\section{Conclusion}

This thesis investigated the construction of a post-quantum threshold batch identity-based encryption (BTE) scheme. The primary motivation for this research stems from the operational requirements of encrypted mempools, where transaction contents must remain entirely concealed until their network ordering is definitively established. Concurrently, decryption authority must be distributed among multiple independent parties to eliminate centralized points of failure. Because existing batch threshold encryption protocols rely predominantly on pairing-based assumptions—which are inherently vulnerable to quantum adversaries—developing a lattice-based alternative represents a critical advancement toward securing decentralized infrastructure in the post-quantum era.

To resolve this, the proposed methodology integrates lattice-based batched decryption techniques with partial lattice trapdoors. Under this architecture, the decryption capability of a central lattice trapdoor is mathematically distributed across multiple authorities. Each participating entity holds only a partial trapdoor, enabling them to locally compute a partial pre-decryption key. Only a predefined, qualified quorum of these authorities can successfully aggregate their respective shares to reconstruct the final batch decryption key. This framework establishes a direct, algebraic approach for distributing decryption capabilities within lattice-based identity-based encryption systems.

Furthermore, this work formally analyzes the correctness and security of the proposed construction. Correctness is derived from the structural properties of the shifted multi-preimage sampler, the underlying sampling procedure studied in~\cite{EC:ABB10}, and the precise reconstruction capabilities of partial lattice trapdoors. From a security perspective, the scheme is proven within a selective security model through a rigorous sequence of hybrid games. These security guarantees are fundamentally anchored in the hardness of established lattice problems, namely Learning with Errors (LWE) and the Short Integer Solution (SIS), augmented by statistical indistinguishability arguments relying on Gaussian sampling and noise flooding techniques.

Ultimately, the proposed scheme represents a foundational theoretical step toward post-quantum threshold encryption for decentralized mempools. While the current construction necessitates further research—particularly regarding concrete parameter optimization, practical implementation efficiency, and the pursuit of tighter security reductions against active adversaries—this work definitively demonstrates that partial lattice trapdoors are a highly promising cryptographic primitive for achieving distributed, post-quantum trust in identity-based infrastructures.

\section{Limitations and Future Work}

While the proposed construction provides a solid and mathematically rigorous foundation, its performance guarantees remain purely asymptotic. Consequently, the scheme in its present form requires further optimization before it can be efficiently deployed in practical production environments. Additionally, the size of the public matrix grows quadratically with the threshold parameter, $\mathcal{O}(t^2)$, which inherently restricts the protocol's viability to scenarios where the number of participating authorities remains relatively small. Furthermore, because the protocol relies entirely on unstructured integer lattices, it suffers from slower arithmetic operations and higher memory consumption. A promising direction for future improvement is to transition the framework to structured lattice variants, specifically Module Learning with Errors~\cite{JACM:LyuPeiReg13} and Ring Learning with Errors~\cite{DCC:LanSte14}

Another structural limitation is the requirement that participating authorities must know the exact composition of the resolving quorum (the approved set) in advance to compute the correct partial decryption shares. This requirement heavily constrains the dynamic flexibility of the network, as the committee cannot be fluidly selected at the time of decryption.

Crucially, the current construction is analyzed under a restricted, highly selective random oracle threat model and remains highly vulnerable to active adversaries. In real-world deployments, adversaries can behave arbitrarily—adaptively manipulating server responses or deliberately injecting malformed shares—and the current scheme is not proven secure against such complex attacks. In the broader landscape of threshold cryptography, mitigating active adversaries typically requires robust countermeasures, such as Verifiable Secret Sharing (VSS) \cite{FOCS:Feldman87} or chosen-ciphertext (CCA) secure threshold frameworks \cite{JC:ShoGen00}, which inherently rely on extensive communication overhead. However, it is important to emphasize that the underlying tool utilized in this proposed construction represents the first fully post-quantum non-interactive key distribution protocol. Consequently, embedding standard active-security measures without destroying this fundamental non-interactive property remains a highly non-trivial open challenge for future research. Moreover, because the random oracle is an idealized model, constructing a post-quantum secure TBIBE protocol in the standard model remains a significant open problem that requires further intensive research.

Finally, to eliminate the key escrow vulnerability and the critical single point of failure inherent in a centralized Private Key Generator (PKG), the architecture must shift toward a distributed framework. To resolve this, future iterations will explore the integration of a Distributed Key Generation (DKG) protocol tailored for identity-based cryptography \cite{SCN:KatGol10}. By leveraging multiparty computation across a network of independent Key Generation Centers (KGCs) or utilizing aggregatable DKG variants suitable for decentralized systems \cite{EC:GKLN21}, the master secret can be generated collaboratively and remain permanently fragmented. This ensures that complete decryption authority is never materialized in a single location, providing the strict trustlessness required for deployment in permissionless blockchain environments.